\begin{figure}
    \centering
    \includegraphics[width=0.35\linewidth]{images/markov.pdf}
    \hspace{0.15\linewidth}
    \includegraphics[width=0.35\linewidth]{images/striated.pdf}
    \begin{flushleft}
    \vspace{-5cm}
    ${\color{cred}\st}$ \hspace{6.75cm} ${\color{cred}\st}$ \\
    \vspace{.04cm}
    ${\color{cblue}\at}$ \hspace{6.75cm} ${\color{cblue}\at}$ \\
    \vspace{.95cm}
    {\color{cred}\Large\vdots} \hspace{6.75cm} \\
        \vspace{-0.95cm}\hspace{7.2cm}
        {\color{cred}\Large\vdots} \\
    \vspace{.1cm}
    {\color{cblue}\Large\vdots}  \hspace{6.75cm} \\
        \vspace{-0.95cm}\hspace{7.2cm}
        {\color{cblue}\Large\vdots} \\
    \vspace{.85cm}
    ${\color{cred}\bs_{t+5}}$  \hspace{6.4cm} ${\color{cred}\bs_{t+5}}$ \\
    \vspace{.02cm}
    ${\color{cblue}\ba_{t+5}}$ \hspace{6.4cm} ${\color{cblue}\ba_{t+5}}$ \\
    \vspace{.05cm}
    \vspace{.2cm}
    \end{flushleft}
    \vspace{.05cm}
    \caption{
    \textbf{(Attention patterns)}
    We observe two distinct types of attention masks during trajectory prediction.
    In the first, both {\color{cred}states} and {\color{cblue}actions} are dependent primarily on the immediately preceding transition, corresponding to a model that has learned the Markov property.
    The second strategy has a striated appearance, with state dimensions depending most strongly on the same dimension of multiple previous timesteps.
    Surprisingly, actions depend more on past actions than they do on past states, reminiscent of the action smoothing used in some trajectory optimization algorithms \citep{nagabandi2019pddm}.
    The above masks are produced by a first- and third-layer attention head during sequence prediction on the hopper benchmark; reward dimensions are omitted for this visualization.$^1$
    }
    \label{fig:attention}
\end{figure}